\centerline{\bf \Huge Теория чиселок}

\begin{flushright}
        \scriptsize{Я очень устал\\
                    Мне хочется спать\\
                    Мне хочется знать...}
\end{flushright}

\problem{1}
Существуют ли такие двузначные числа $\overline{a b}$ и $\overline{c d}$, что $\overline{a b} \cdot \overline{c d}=\overline{a b c d}$ ?

\problem{2}
При каких натуральных $n$ число $n^3+2 n^2+11$ является кубом натурального числа?

\problem{3}
О натуральных числах $a, p, q$ известно, что $a p+1$ делится на $q$, а $a q+1$ делится на $p$. Докажите, что $a>\frac{p q}{2(p+q)}$.

\problem{4}
 Пусть $a$, $b$, $c$ --- произвольные натуральные числа. Докажите, что числа $a^2+b+c$, $b^2+c+a$, $c^2+a+b$ не могут одновременно быть точными квадратами.

\problem{5}
Докажите, что из любых десяти последовательных натуральных чисел можно выбрать число, взаимно простое с остальными девятью.\\
Докажите, что произведение десяти последовательных натуральных чисел не может быть точным квадратом.

\problem{6}
 Вася выписал в строку 100 последовательных чисел. Во второй строке под каждым числом он написал его собственный делитель. В третьей строке он под каждым числом из второй строки написал его собственный делитель и т.\,д., пока не получилось 2022 строк. Могут ли в каждой строке стоять последовательные числа?

\problem{7}
  Пусть $p > 5$ --- простое число и $S:=\{p - n^2 \colon n \in \mathbb{N}, \, n^2 < p\}$. Докажите, что во множестве~$S$ найдутся два числа, одно из которых делит другое

\problem{8}
 Найти все тройки натуральных чисел $(a,b,c)$, такие, что $2^a+2^b+1$ делится на $2^c+1$.

 